% Compile with pdflatex raul-kolaric-en.tex (or tectonic).
\documentclass[10pt,letterpaper]{article}
\usepackage[letterpaper,margin=0.45in]{geometry}
\usepackage[T1]{fontenc}
\usepackage[utf8]{inputenc}
\usepackage{mathptmx}
\usepackage{enumitem}
\usepackage[hidelinks]{hyperref}
\pagestyle{empty}
\linespread{0.95}
\setlength{\parindent}{0pt}
\setlength{\parskip}{0pt}
\setlist[itemize]{leftmargin=1.3em,itemsep=1pt,parsep=0pt,topsep=2pt}
\newcommand{\ressection}[1]{\par\vspace{6pt}{\large\bfseries #1}\par\vspace{2pt}\hrule\vspace{3pt}}
\newcommand{\role}[4]{\textbf{#1}\hfill #2\par\textbf{#3}\hfill #4\par}
\begin{document}
\begin{center}
{\LARGE\bfseries RAUL KOLARIC}\\[3pt]
{\small São Paulo, Brazil $\bullet$
\href{mailto:rlkolaric@gmail.com}{\underline{rlkolaric@gmail.com}} $\bullet$
\href{tel:+5511993462745}{\underline{+55 11 99346-2745}} $\bullet$
\href{https://www.linkedin.com/in/raulkolaric/}{\underline{LinkedIn}} $\bullet$
\href{https://github.com/raulkolaric}{\underline{GitHub}} $\bullet$
\href{https://rkolaric.com/}{\underline{Portfolio}}}
\end{center}
\vspace{-7pt}
\ressection{Education}
\textbf{Pontifícia Universidade Católica de São Paulo (PUC-SP)}\hfill 2025 -- 2028 (Expected)
\begin{itemize}
\item Bachelor of Science in Computer Science.
\end{itemize}
\ressection{Experience}
\role{Microsoft}{São Paulo, Brazil}{Business Program Management Intern}{September 2026 -- Present}
Supporting cross-functional business and technical initiatives through data analysis, program optimization, and AI-enabled solutions.
\par\vspace{8pt}
\role{\href{https://upp.com.br/}{Up.p}}{São Paulo, Brazil}{Product Management Intern}{May 2026 -- July 2026}
As an intern at Up.p, a fintech leader in payroll-deducted financial products with over \$200 million in total volume, I contributed to the development and optimization of high-scale financial products.
\begin{itemize}
\item Led the prototyping and technical documentation for the new Employer Portal (Consignado CLT), translating complex business requirements into actionable product features for the development team.
\item Streamlined internal workflows by automating key operational processes, resulting in improved team productivity and significant time savings.
\item Managed performance optimizations and content updates for the company's main website using WordPress, ensuring high availability and consistent user engagement.
\item Authored technical specifications and business requirements documentation, bridging the gap between leadership needs and engineering implementation.
\end{itemize}
\vspace{8pt}
\role{\href{https://mapari.app/}{Mapari}}{São Paulo, Brazil}{Cofounder \& Full Stack Developer}{January 2024 -- February 2026}
I participated in the development of the Mapari application, focused on mapping the real-time movement of bars and events in the city of São Paulo.
\begin{itemize}
\item Developed an administrative control panel in Next.js, React, and TypeScript, optimized for the real-time visualization of data from over 80,000 daily requests.
\item Executed the front-end implementation of a landing page based on the Figma design, using Next.js and TypeScript. The architecture was based on modular componentization and integrated market libraries to ensure clean, reusable, and high-performance code.
\item Modeled and implemented a relational data architecture in PostgreSQL (Supabase), structuring relationships between multiple tables and applying industry normalization standards to ensure integrity and performance under millions of monthly requests.
\item Built the backend layer through REST APIs in Next.js/TypeScript, leveraging the framework's native route structure for performant and type-safe endpoints.
\item Implemented an infrastructure solution on AWS, using EC2 with cron jobs for scheduled tasks and S3 for scalable file storage.
\end{itemize}
\ressection{Leadership \& Volunteering}
\role{PUC-TECH}{Academic League}{Member}{June 2025 -- Present}
\begin{itemize}
\item Led a team of four people in the development of a platform for identifying accidents on São Paulo state highways. Through machine learning, the project aimed to determine accident mitigation points based on official state data, gaining valuable insights into leadership, people management, and the project structure and lifecycle.
\item Drove the successful and timely development of the platform by establishing organizational routines (meetings and checkpoints), managing deadlines and deliverables, and overseeing the fidelity of the design implementation.
\item Collaborated on the organizing committee for PUC-SP's academic computer science week, assisting speakers and managing the event schedule.
\end{itemize}
\vspace{8pt}
\role{INSTITUTO FEFIG}{NGO}{Volunteer}{October 2025}
\begin{itemize}
\item Operational support in fundraising for inter-state educational projects.
\end{itemize}
\ressection{Technical Skills}
\textbf{Spoken Languages:} Portuguese (Native), English (Fluent, C1/C2), Spanish (Intermediate).\\
\textbf{Languages:} TypeScript, JavaScript, Python, C, HTML/CSS.\\
\textbf{Frameworks \& Web:} Next.js, React, TailwindCSS, FastAPI, REST APIs.\\
\textbf{Data \& AI:} Pandas, Scikit-learn, Matplotlib, BeautifulSoup.\\
\textbf{Automation \& DevOps:} Playwright, Docker, AWS (S3, EC2, RDS), Git/GitHub, Linux.\\
\textbf{Databases:} PostgreSQL, Supabase, Relational Modeling (SQL)\\
\textbf{Design \& Tools:} Figma, Canva, Agile (Scrum).
\end{document}
